\documentclass[11pt]{article}
\usepackage{fltcourse}

\title{Level 6: From $\ell$ to $3$}
\author{Kevin Buzzard}
\date{}

\begin{document}
\maketitle

\noindent
In the previous level we reduced Fermat's Last Theorem to Statement $B_5$: for
$\ell\geq5$ prime, every continuous hardly ramified representation
$\rho\colon\GQ\to\Aut_{\F_\ell}(V)$ on a two-dimensional $\F_\ell$-vector space
is reducible. All the terminology --- coefficient ring, hardly ramified,
cyclotomic character $\chi$, flat at $\ell$ --- is as in the previous level;
recall in particular that the definition of ``hardly ramified'' makes sense
over any coefficient ring of odd residue characteristic, not only $\F_\ell$.

\medskip\noindent
Wiles' great insight was to study the prime $3$. We too shall pass from the
prime $\ell$ to the prime $3$. But the Frey curve is long gone: our
representations are now free-standing objects, and there is no elliptic curve
to reduce modulo $3$. Passing from $\ell$ to $3$ must therefore be done
entirely on the level of Galois representations. This level introduces the
three deep statements that make the passage possible, and shows that together
they imply $B_5$ --- and hence FLT. Proving the three statements themselves is
the business of every subsequent level.

\section{The three heads}

We split $B_6$ into three statements $B_{6a}$, $B_{6b}$, $B_{6c}$, and set
$B_6:=B_{6a}\wedge B_{6b}\wedge B_{6c}$. In all three, ``representation'' means
``continuous representation on a free rank-two module over a coefficient ring
of odd residue characteristic'', and ``hardly ramified'' is as before.

\begin{Bn}{B_{6a}}
\textup{(Lifting.)} Let $\ell\geq5$ be prime and let
$\rhobar\colon\GQ\to\Aut_{\F_\ell}(V)$ be an irreducible hardly ramified
$\F_\ell$-representation. Then there exist a finite extension $K/\Ql$ with ring
of integers $\calO$ and maximal ideal $\m$, a free $\calO$-module $\calV$ of
rank $2$, and a hardly ramified representation $\rho\colon\GQ\to\Aut_\calO(\calV)$
whose reduction $\rho\bmod\m$ is isomorphic to the base change of $\rhobar$ to
$\calO/\m$.
\end{Bn}

\noindent
In slogan form: \emph{an irreducible hardly ramified mod $\ell$ representation
lifts to an $\ell$-adic one}. (The reader who is still thinking about the Frey
curve should note that its $\ell$-adic Tate module will not do as a lift: it is
ramified at every prime dividing $abc$, so is not hardly ramified.)

\begin{Bn}{B_{6b}}
\textup{(Spreading out.)} Let $\ell\geq5$ be prime, let $K/\Ql$ be a finite
extension with ring of integers $\calO$ and maximal ideal $\m$, let $\calV$ be
a free rank-two $\calO$-module, and let $\rho\colon\GQ\to\Aut_\calO(\calV)$ be a
hardly ramified representation. Assume that the reduction $\rho\bmod\m$ is
absolutely irreducible and can be conjugated into $\GL_2(\F_\ell)$ \textup{(}i.e.\ is
the base change of an irreducible $\F_\ell$-representation\textup{)}. Then:
\begin{enumerate}
\item[\textup{(i)}] there is a number field $M\subseteq K$ and elements
$t_p\in M$, for all primes $p\neq2,\ell$, with $\tr\rho(\Frob_p)=t_p$; and
\item[\textup{(ii)}] for every maximal ideal $P$ of the ring of integers
$\calO_M$ of residue characteristic $p\neq2$, there is a hardly ramified
$\calO_{M,P}$-representation $\rho_P\colon\GQ\to\GL_2(\calO_{M,P})$ with
$\tr\rho_P(\Frob_q)=t_q$ for all primes $q\neq 2,p,\ell$.
\end{enumerate}
\end{Bn}

\noindent
In slogan form: \emph{an $\ell$-adic hardly ramified representation whose
residual representation descends to $\F_\ell$ spreads out into a compatible
family of hardly ramified $p$-adic representations, one for each odd residue
characteristic $p$}. The traces $t_p$ live in a single number field $M$ and are
common to the whole family: this is what ``compatible'' means. The remarkable
feature is that there is \emph{no} hypothesis that anything be modular.

\begin{Bn}{B_{6c}}
\textup{(Classification at $3$.)} Let $K/\Q_3$ be a finite extension with ring
of integers $\calO$, let $\calV$ be a free rank-two $\calO$-module, and let
$\rho\colon\GQ\to\Aut_\calO(\calV)$ be a hardly ramified representation. Then,
for a suitable basis,
\[
  \rho\;\cong\;\begin{pmatrix}\chi & * \\ 0 & 1\end{pmatrix},
\]
where $\chi$ is the $3$-adic cyclotomic character. In particular
$\tr\rho(\Frob_p)=1+p$ for every prime $p\neq2,3$.
\end{Bn}

\begin{remark}[compatible families]
It is worth pausing over what $B_{6b}$ asserts. We learn as undergraduates that
two finite-dimensional representations with the same character are isomorphic
--- but that is a statement about \emph{one} representation of \emph{one} group.
The family $(\rho_P)_P$ consists of genuinely different representations of
$\GQ$, with values in different rings, and yet their traces of Frobenius agree
and lie in a common number field. The simplest instance is already striking:
the $2$-adic and $3$-adic cyclotomic characters $\chi_2\colon\GQ\to\Z_2^\times$
and $\chi_3\colon\GQ\to\Z_3^\times$ have disjoint kernels (cutting out the
disjoint fields $\Q(\mu_{2^\infty})$ and $\Q(\mu_{3^\infty})$) and are in every
sense different representations; yet for every prime $p\geq5$ one has
$\chi_2(\Frob_p)=p=\chi_3(\Frob_p)$, a common rational integer. This
coincidence of traces across a family is the first shadow of the theory of
motives, and it is exactly the mechanism by which we shall transport
information from $\ell$ to $3$.
\end{remark}

\section{An elementary lemma}

The reduction of $B_5$ to $B_6$ passes through $B_{6b}$, whose hypothesis
demands \emph{absolute} irreducibility of the residual representation. The
following lemma, which uses nothing beyond the cyclotomic determinant and the
existence of complex conjugation, closes the gap between irreducibility and
absolute irreducibility for our representations.

\begin{lemma}\label{lem:absirr}
Let $\ell$ be an odd prime and let $\rhobar\colon\GQ\to\GL_2(\F_\ell)$ be an
irreducible representation with $\det\rhobar=\chibar$, the mod $\ell$
cyclotomic character. Then $\rhobar$ is absolutely irreducible, i.e.\ remains
irreducible after extension of scalars to $\overline{\F}_\ell$.
\end{lemma}

\begin{proof}
Let $c\in\GQ$ be a complex conjugation (the restriction to $\Qbar\subseteq\C$
of complex conjugation for some embedding). Then $c^2=1$ and $c$ inverts every
root of unity, so $\chibar(c)=-1$. Hence $\rhobar(c)$ is an involution of
determinant $-1$; as $\ell$ is odd its eigenvalues are the distinct elements
$+1,-1$ of $\F_\ell$, so $\rhobar(c)$ is diagonalisable with both eigenlines
defined over $\F_\ell$. Now suppose $\overline{\F}_\ell^{\,2}$ had a
$\rhobar(\GQ)$-stable line $\Lambda$. Then $\Lambda$ is $\rhobar(c)$-stable,
hence equals one of the two eigenlines of $\rhobar(c)$, which are defined over
$\F_\ell$. A $\GQ$-stable line defined over $\F_\ell$ contradicts the
irreducibility of $\rhobar$ over $\F_\ell$. So no such $\Lambda$ exists and
$\rhobar$ is absolutely irreducible.
\end{proof}

\section{The boss: $B_6$ implies FLT}

\begin{theorem}[Boss of Level 6]\label{thm:boss6}
Statements $B_{6a}$, $B_{6b}$ and $B_{6c}$ together imply Fermat's Last
Theorem, modulo results known in the 1980s.
\end{theorem}

\begin{proof}
By the boss theorem of the previous level it suffices to prove $B_5$. We argue
by contradiction: suppose $\ell\geq5$ is prime and
$\rhobar\colon\GQ\to\Aut_{\F_\ell}(V)$ is an \emph{irreducible} hardly ramified
$\F_\ell$-representation, and let us derive a contradiction. (This proves that
no such $\rho$ exists, i.e.\ that every hardly ramified
$\F_\ell$-representation is reducible, which is $B_5$.)

\medskip\noindent\emph{Lift.} By $B_{6a}$ there is a finite extension $K/\Ql$
with ring of integers $\calO$ and maximal ideal $\m$, and a hardly ramified
representation $\rho\colon\GQ\to\Aut_\calO(\calV)$ with
$\rho\bmod\m\cong\rhobar\otimes_{\F_\ell}(\calO/\m)$.

\medskip\noindent\emph{Verify the hypothesis of $B_{6b}$.} By
Lemma~\ref{lem:absirr}, $\rhobar$ is absolutely irreducible; hence so is its
base change $\rho\bmod\m$, which moreover is manifestly conjugate into
$\GL_2(\F_\ell)$, being the base change of the $\F_\ell$-representation
$\rhobar$. So $B_{6b}$ applies.

\medskip\noindent\emph{Spread out to $3$.} By $B_{6b}$ there is a number field
$M\subseteq K$, elements $t_p\in M$ with $\tr\rho(\Frob_p)=t_p$ for $p\neq2,\ell$,
and a compatible family $(\rho_P)$. Since $\ell\geq5$, the prime $3$ satisfies
$3\neq2,\ell$; choose a maximal ideal $P$ of $\calO_M$ above $3$, of residue
characteristic $3$. Then $\rho_P\colon\GQ\to\GL_2(\calO_{M,P})$ is a hardly
ramified $3$-adic representation with $\tr\rho_P(\Frob_q)=t_q$ for all
$q\neq2,3,\ell$.

\medskip\noindent\emph{Classify at $3$.} By $B_{6c}$, $\rho_P$ is conjugate to
$\left(\begin{smallmatrix}\chi_3&*\\0&1\end{smallmatrix}\right)$, so
$\tr\rho_P(\Frob_q)=\chi_3(q)+1=q+1$ for every $q\neq2,3$. Combining with the
previous paragraph,
\[
  t_q=q+1\qquad\text{for all primes }q\neq2,3,\ell.
\]
Reducing modulo $\m$ (recall $t_q=\tr\rho(\Frob_q)$ and
$\rho\bmod\m=\rhobar\otimes k$), we obtain
\[
  \tr\rhobar(\Frob_q)=q+1\qquad\text{for all primes }q\neq2,3,\ell,
\]
an identity in $\F_\ell$. But $q+1=\chibar(\Frob_q)+1=\tr\bigl((\mathbf1\oplus\chibar)(\Frob_q)\bigr)$,
so $\rhobar$ and $\mathbf1\oplus\chibar$ have equal traces on every $\Frob_q$
with $q\neq2,3,\ell$.

\medskip\noindent\emph{Chebotarev and Brauer--Nesbitt.} Both $\rhobar$ and
$\mathbf1\oplus\chibar$ have finite image and are unramified outside $\{2,\ell\}$;
let $L/\Q$ be a finite Galois extension through which both factor, so both are
inflated from representations of $\Gal(L/\Q)$, and $L$ is unramified outside
$\{2,\ell\}$. Every prime $q\neq2,3,\ell$ is then unramified in $L$ and supplies
a Frobenius conjugacy class $\Frob_q\subseteq\Gal(L/\Q)$. By the Chebotarev
density theorem (1920s, certainly available), every conjugacy class of
$\Gal(L/\Q)$ arises as $\Frob_q$ for infinitely many such $q$. As trace is a
class function, we conclude that $\rhobar$ and $\mathbf1\oplus\chibar$ have
\emph{equal character} on all of $\Gal(L/\Q)$. Both are semisimple ($\rhobar$
because it is irreducible), so by the Brauer--Nesbitt theorem a semisimple
representation is determined up to isomorphism by its character. (In residue
characteristic $\ell$ one needs equal characteristic polynomials, not merely
equal traces; but both representations have determinant $\chibar$ and $\ell>2$,
so equal traces do give equal characteristic polynomials.) Hence
$\rhobar\cong\mathbf1\oplus\chibar$.

\medskip\noindent
But $\mathbf1\oplus\chibar$ is visibly reducible, contradicting the
irreducibility of $\rhobar$. This contradiction proves $B_5$, and with it the
theorem.
\end{proof}

\begin{remark}
The moral of the level is that profound theorems about $p$-adic Galois
representations imply Fermat's Last Theorem, and that the passage from the
prime $\ell$ to the manageable prime $3$ is effected purely by keeping track
of traces of Frobenius in a compatible family. What remains is to prove
$B_{6a}$, $B_{6b}$ and $B_{6c}$. Of the three, $B_{6c}$ is by far the most
accessible: its proof, in the next level, rests on a classification of mod $3$
hardly ramified representations by way of discriminant bounds, and needs
nothing from after the 1980s. Statements $B_{6a}$ and $B_{6b}$ are the two
faces of a single automorphy lifting theorem, and occupy the later levels.
\end{remark}

\end{document}
